\chapter*{คำนำ}
\addcontentsline{toc}{chapter}{คำนำ}

วิกฤตการณ์การเปลี่ยนแปลงสภาพภูมิอากาศ (Climate Change) และมลพิษทางอากาศจากฝุ่นละอองขนาดเล็กไม่เกิน 2.5 ไมครอน (PM2.5) นับเป็นภัยคุกคามทางสิ่งแวดล้อมและสาธารณสุขที่รุนแรงที่สุดในศตวรรษที่ 21 ตำราวิชาการเรื่อง \textbf{“ฟิสิกส์สิ่งแวดล้อม การตรวจวัด PM2.5 และนวัตกรรมพลังงานชีวมวล” (Environmental Physics, PM2.5 Deep Learning Sensing, and Biomass Clean Energy Innovations)} เล่มนี้ จัดทำขึ้นเพื่อบูรณาการองค์ความรู้ทางฟิสิกส์ประยุกต์ ทัศนศาสตร์บรรยากาศ ปัญญาประดิษฐ์เพื่อการตรวจวัดมลพิษ และวิศวกรรมพลังงานความร้อนสะอาดเข้าด้วยกันอย่างเป็นระบบ

เนื้อหาในส่วนแรกของตำรามุ่งเน้นฟิสิกส์ของชั้นบรรยากาศ พลศาสตร์การแพร่กระจายของอนุภาคแขวนลอย ทฤษฎีการกระเจิงแสงของอนุภาคฝุ่นแบบมี (Mie Scattering) และนวัตกรรมการตรวจวัดค่าความเข้มข้นของฝุ่น PM2.5 ด้วยภาพถ่ายสมาร์ทโฟนร่วมกับโครงข่ายประสาทเทียมแบบน้ำหนักเบา (Mobile Deep Learning) ซึ่งเปิดโอกาสให้ประชาชนและชุมชนสามารถเข้าถึงการเฝ้าระวังคุณภาพอากาศได้ด้วยตนเอง

ในส่วนที่สองนำเสนอเทอร์โมไดนามิกส์ของการเผาไหม้และฟิสิกส์พลังงานชีวมวล โดยบูรณาการผลงานวิจัยจริงของผู้เขียนเกี่ยวกับการพัฒนาถ่านอัดแท่งเสริมประสิทธิภาพความร้อนด้วยคาร์บอนแบล็ค (Carbon Black Briquettes) การเปลี่ยนของเสียทางการเกษตร เช่น เปลือกมะม่วงหิมพานต์และเศษชีวมวลเป็นพลังงานสะอาด การผลิตน้ำมันชีวภาพและถ่านชีวภาพผ่านกระบวนการไพโรไลซิส ตลอดจนการประเมินคาร์บอนเครดิตตามโครงการลดก๊าซเรือนกระจกภาคสมัครใจของประเทศไทย (T-VER) ทั้งด้านพลังงานและป่าไม้

ผู้เขียนหวังว่าตำราเล่มนี้จะเป็นรากฐานทางวิชาการที่สำคัญแก่นักศึกษา นักวิจัย และผู้กำหนดนโยบายในการขับเคลื่อนประเทศไทยสู่สังคมคาร์บอนต่ำและสิ่งแวดล้อมที่ยั่งยืน

\vspace{1.5cm}
\begin{flushright}
    \textbf{ผู้ช่วยศาสตราจารย์ ดร.ชีวะ ทัศนา}\\[0.2cm]
    สาขาวิชาฟิสิกส์ คณะวิทยาศาสตร์และเทคโนโลยี\\[0.1cm]
    มหาวิทยาลัยราชภัฏรำไพพรรณี\\[0.1cm]
    พุทธศักราช 2569
\end{flushright}
